The Serial Peripheral Interface (SPI) peripheral is a high speed, full-duplex, synchronous, multipoint serial interface. It transfers one bit at a time. The SPI bus requires three wires (SCK, MOSI, and MISO), along with one CS wire for every slave. The main features of the SPI peripheral are listed below:

\begin{itemize}
	\item Master or slave mode (except \peripheral{SPI0}, which is master only)
	\item LSB-first or MSB-first transfers
	\item 8-, 16-, or 32-bit transfer lengths
	\item Master mode supports SPI modes 0-3
	\item Slave mode supports SPI modes 1 and 3
	\item Optional byte swap for multi-byte transfers with sytems of different endiannesses
	\item Programmable transfer rate
	\item Continuous transfer operation to reduce dead time
	\item Interrupts for transfer complete and TX register empty
	\item Native memory read access for SPI flash chips (on \peripheral{SPI0} only)
\end{itemize}


\subsection{SPI Pins}

\begin{figure}[htpb]
	\includesvg{figures/spi-connection-diagram}
	\caption{SPI Connection Diagram}
	\label{fig:spi-connection-diagram}
\end{figure}

The SPI bus has four pins: \pin{SCKx}, \pin{MOSIx}, \pin{MISOx}, and \pin{CSx} (not present on \peripheral{SPI0}). A connection diagram for the SPI bus is shown in Figure \ref{fig:spi-connection-diagram}.

\pin{SCKx} is the serial clock pin, which governs when each data bit is both latched and updated. It must be attached to the SCK pin on every other device that shares the SPI bus. When the corresponding \register{PxSEL[y]} bit is a 1, the \peripheral{SPIx} peripheral seizes the \texttt{DIR}, \texttt{OUT}, and \texttt{REN} controls of \pin{SCKx}.

\pin{MOSIx} is the master out slave in pin, which transfers each data bit from the master to the slave. It is an output when the SPI port is enabled and in master mode, otherwise it is a high impedance input. It must be attached to the MOSI pin on every other device that shares the SPI bus. When the corresponding \register{PxSEL[y]} bit is a 1, the \peripheral{SPIx} peripheral seizes the \texttt{DIR}, \texttt{OUT}, and \texttt{REN} controls of \pin{MOSIx}.

\pin{MISOx} is the master in slave out pin, which transfers each data bit from the slave to the master. It is an output when the SPI port is enabled, in slave mode, and the \pin{CSx} pin is asserted (low), otherwise it is a high impedance input. It must be attached to the MISO pin on every other device that shares the SPI bus. When the corresponding \register{PxSEL[y]} bit is a 1, the \peripheral{SPIx} peripheral seizes the \texttt{DIR}, \texttt{OUT}, and \texttt{REN} controls of \pin{MISOx}.

\pin{CSx} is the chip select pin for the \peripheral{SPIx} peripheral on this chip. It is a high impedance input that waits to be asserted (pulled low) by an external master if this \peripheral{SPIx} peripheral is in slave mode, which begins the SPI transfer. One CS wire needs to be connected between a GPIO pin of the master and the CS pin of the slave for every slave on the bus. Each slave CS wire must be attached to a dedicated GPIO pin on the master in a one-to-one fashion. The \peripheral{SPIx} peripheral never seizes control of the \pin{CSx} pin (regardless of the state of the corresponding \register{PxSEL[y]} bit), and as such it must be configured as a high impedance input manually.

\subsection{Generalized SPI Transfer Process}

\begin{figure}[htpb]
	\begin{tikztimingtable}[timing/font=\rmfamily, timing/d/text/.append style={font=\rmfamily}, xscale=1.5, timing/slope=0.2, timing/dslope=0.2, timing/d/background/.style={fill=white}]
		CPOL = 0 & LL 15{T} LL \\
		CPOL = 1 & HH 15{T} HH \\
		CS & H 17L H     \\
		\\
		Cycle \# & U     R 8{2Q} 2U    \\
		MISO & D{z}  R 8{2Q} 2D{z} \\
		MOSI & D{z}  R 8{2Q} 2D{z} \\
		\\
		Cycle \# & UU    R 8{2Q} U    \\
		MISO & D{z}U R 8{2Q} D{z} \\
		MOSI & D{z}U R 8{2Q} D{z} \\
		\extracode
		% Add vertical lines in two colors
		\begin{pgfonlayer}{background}
			\begin{scope}[semitransparent,semithick]
				\vertlines[red]{2.1,4.1,...,17.1}
				\vertlines[blue]{3.1,5.1,...,17.1}
			\end{scope}
		\end{pgfonlayer}
		% Add big group labels
		\begin{scope}
			[shift={(-4em,-0.5)},anchor=east]
			\node at (  0, 0) {SCK};
			\node at (1ex,-9) {CPHA = 0};
			\node at (1ex,-17) {CPHA = 1};
		\end{scope}
	\end{tikztimingtable}
	\caption{SPI Timing Diagram}
	\label{fig:spi-timing-diagram}
\end{figure}

When no transfer is occurring, the SPI bus is idle. The transfer begins when the master asserts (pulls low) the CS wire attached to the desired slave. Immediately, the master seizes control of the MOSI wire by making its MOSI pin an output, and the slave does the same for the MISO wire. Note that the master always has control of the SCK and CS wires, which are configured as outputs. Then, the master issues SCK clock pulses in increments of 8, 16, or 32 (one frame of data). On every transition of SCK, both the master and slave alternate between updating and latching the values of MOSI and MISO, the master updating MOSI or latching MISO, and the slave updating MISO or latching MOSI. When the wire is updated, the next bit of data is pushed out of the TX shift register and written to the wire. When the wire is latched, that bit is read from the wire and pushed into the RX shift register. As the SPI port is full-duplex, the master and slave both send data to the other simultaneously. After each frame of data has been transferred, both the master and slave RX shift registers contain their respective received words (the 8-, 16-, or 32-bit data transmitted in the frame). At that point, the master may either continue the process by transferring another frame, or cease the transfer by deasserting (pulling high) the CS wire. The frames that the master issues while the slave's CS wire is low is called a packet, and the master may issue as many frames in a packet as it wants before deasserting CS.

The order of the update/latch process, as well as the polarity of the SCK wire, are determined by two 1-bit parameters: clock polarity (CPOL) and clock phase (CPHA). CPOL determines the idle state and transition direction of SCK, while CPHA determines when the data in MOSI and MISO is updated and latched. These parameters together are referred to as the SPI mode, or SPIMODE, shown in Table \ref{t:SPIMODE-key}. Note that this SPI peripheral supports all four SPI modes while in master mode, but only supports SPI modes 1 and 3 while in slave mode (CPHA = 1). Figure \ref{fig:spi-timing-diagram} shows a SPI timing diagram for all four SPI modes, with a single-frame packet with 8-bit frames.

\begin{table}[htb]
	\centering
	\begin{tabular}{ c c c c c c c }
	\caption{SPIMODE Key} \label{t:SPIMODE-key}
	\textbf{SPIMODE} & \textbf{CPOL} & \textbf{CPHA} & \textbf{SCK Idle} & \textbf{Data Update Edge} & \textbf{Data Latch Edge} \\ \hline
	0 & 0 & 0 & Low  & Trailing/Falling & Leading/Rising \\
	\rowcolor{tablehighlightcolor}
	1 & 0 & 1 & Low  & Leading/Rising & Trailing/Falling \\
	2 & 1 & 0 & High  & Trailing/Rising & Leading/Falling \\
	\rowcolor{tablehighlightcolor}
	3 & 1 & 1 & High  & Leading/Falling & Trailing/Rising \\
	\hline
	\end{tabular}
\end{table}

Note that during the very first frame of a transfer, the slave is required to send data to the master on the MISO wire, regardless of if it has any valid data to send. If it has no data to send, a slave will issue a ``dummy'' word, which the master must ignore.



\subsection{Bit and Byte Ordering}

The SPI peripheral has the ability to transfer the most significant bit (MSB) first or the least significant bit (LSB) first using \bitfield{SPIMSB}. When \bitfield{SPIMSB} is 0, the LSB is transferred first, and when \bitfield{SPIMSB} is 1, the MSB is transferred first.

For 16- and 32-bit transfers, the byte ordering may optionally be swapped using \bitfield{SPITXSB} and \bitfield{SPIRXSB}, which swap the bytes being transmitted and received, respectively. For example, during a 16-bit transfer, if \bitfield{SPITXSB} is enabled, then byte 1 of \register{SPITX} will be transmitted followed by byte 0. If \bitfield{SPIRXSB} is enabled, then the first byte received will be placed in byte 1 of \register{SPIRX}, and the next 8 bits received will be placed in byte 0. In a 32-bit transfer, the byte order would thus be 3, 2, 1, 0. Figure \ref{fig:spi-byte-ordering} illustrates this ordering.

\begin{figure}[htpb]
	\begin{tikztimingtable}[timing/font=\rmfamily, timing/d/text/.append style={font=\rmfamily}, xscale=4, timing/slope=0.1, timing/dslope=0.2]
		No byte swap                & [X] D{byte 0} D{byte 1} D{byte 2} D{byte 3} D{byte 4} D{byte 5} D{byte 6} D{byte 7} D{...} \\
		16-bit transfers, byte swap & [X] D{byte 1} D{byte 0} D{byte 3} D{byte 2} D{byte 5} D{byte 4} D{byte 7} D{byte 6} D{...} \\
		32-bit transfers, byte swap & [X] D{byte 3} D{byte 2} D{byte 1} D{byte 0} D{byte 7} D{byte 6} D{byte 5} D{byte 4} D{...} \\
	\end{tikztimingtable}
	\caption{SPI Byte Ordering}
	\label{fig:spi-byte-ordering}
\end{figure}

Note that the bits in the frame are ordered MSB-first or LSB-first \textit{before} the byte swap operation takes place. As such, a 16-bit, MSB-first transfer with byte swapping activated would transmit the bits 7 down to 0, and then 15 down to 8. Figure \ref{fig:spi-bit-ordering} shows the bit ordering for a 16-bit transfer.

\begin{figure}[htpb]
	\begin{tikztimingtable}[timing/font=\rmfamily, timing/d/text/.append style={font=\rmfamily}, xscale=2, timing/slope=0.1, timing/dslope=0.2]
		No byte swap, LSB-first & [X] D{0} D{1} D{2} D{3} D{4} D{5} D{6} D{7} D{8} D{9} D{10} D{11} D{12} D{13} D{14} D{15} D{...} \\
		Byte swap, LSB-first    & [X] D{8} D{9} D{10} D{11} D{12} D{13} D{14} D{15} D{0} D{1} D{2} D{3} D{4} D{5} D{6} D{7} D{...} \\
		No byte swap, MSB-first & [X] D{15} D{14} D{13} D{12} D{11} D{10} D{9} D{8} D{7} D{6} D{5} D{4} D{3} D{2} D{1} D{0} D{...} \\
		Byte swap, MSB-first    & [X] D{7} D{6} D{5} D{4} D{3} D{2} D{1} D{0} D{15} D{14} D{13} D{12} D{11} D{10} D{9} D{8} D{...} \\
	\end{tikztimingtable}
	\caption{SPI Bit Ordering (16-bit transfer)}
	\label{fig:spi-bit-ordering}
\end{figure}



\subsection{SPI Baud Clock}

When in master mode, the frequency of the \pin{SCKx} clock pin is determined by the frequency of SMCLK and the SPI baud control \bitfield{SPIBR}, given by the expression:
\begin{equation*}
f_\textrm{SCK} = \frac{f_\textrm{SMCLK}}{2 \times (1 + \texttt{SPIBR})}
\end{equation*}

Note that the maximum frequency of \pin{SCKx} is half that of SMCLK.


\subsection{Transmit Buffer Register Queue}

Writing to the SPI transmit buffer register \register{SPIxTX} queues up the word written to it as the next frame to transmit. If in master mode and a transfer is not currently ongoing (i.e., \pin{SCKx} is not clocking), the SPI peripheral immediately begins transmitting the word written to \register{SPIxTX}, and \bitfield{SPITEIF} is set immediately. If a transfer is currently ongoing, the SPI peripheral waits for the current frame to finish transferring. As soon as it finishes, the SPI peripheral begins transmitting the new frame written to \register{SPIxTX} and sets \bitfield{SPITEIF}. If in slave mode, the SPI peripheral latches and begins transferring the last word written to \register{SPIxTX} and sets \bitfield{SPITEIF} whenever the master starts the first clock cycle of \pin{SCKx}.

The TX queue allows the SPI peripheral to constantly transfer data without any dead time in between frames, ensuring maximum transfer speed. However, the user must take care not to overflow the queue, which happens if the \register{SPIxTX} is written two or more times during a single frame transfer. Queue overflows can be prevented by waiting for the \bitfield{SPITEIF} flag to be 1 or \bitfield{SPIBUSY} to be 0 before writing to \register{SPIxTX}.


\subsection{Flags and Interrupts}

The SPI peripheral has two interrupt flags (\bitfield{SPITCIF} and \bitfield{SPITEIF}) and one indicator (\bitfield{SPIBUSY}). If the corresponding interrupts are enabled in the \register{SPIxCR} register and the \peripheral{SPIx} ISR is not masked, then the two interrupt flags will generate an interrupt whenever they are equal to 1. Both interrupt flags can be cleared by writing a 0 to them, which must be done before the \peripheral{SPIx} ISR returns, else the ISR will mistakenly re-execute immediately upon the prior return.

The \bitfield{SPITCIF} interrupt flag is set the moment a SPI transfer is completed, and can only be cleared by writing a 0 to it.

The \bitfield{SPITEIF} interrupt flag is set the moment the SPI TX transmit buffer register \register{SPIxTX} becomes empty (i.e., is ready to queue up the next word to be transmitted). In master mode, if a frame is not currently being transferred, it is set immediately when a new word is written to \register{SPIxTX}. If a frame is currently being transferred, it is set as soon as the current transfer is completed, whereupon the value in \register{SPIxTX} begins to be transmitted. In slave mode, \bitfield{SPITEIF} is set when the master starts the first clock cycle of \pin{SCKx}.

The \bitfield{SPIBUSY} indicator simply indicates if the SPI is busy or not. In master mode, it is 1 if a frame is currently being transferred (i.e. \pin{SCKx} is actively clocking), and 0 if not. Note that \bitfield{SPIBUSY} can therefore be 0 in between frame transfers if a new word to transmit has not been queued before the prior frame finishes transferring. In slave mode, \bitfield{SPIBUSY} is 1 if the \pin{CSx} pin is asserted (driven low), and is 0 if not.

\subsection{Master Mode Transfer Procedure}

To use the SPI peripheral in master mode, the user must first initialize the \peripheral{SPIx} peripheral by configuring the \register{SPIxCR} register. Of note, \bitfield{SPIEN} must be 1 and \bitfield{SPISM} must be 0 to enable the SPI peripheral in master mode. Then, the GPIO pins multiplexed with \pin{SCKx}, \pin{MISOx}, and \pin{MOSIx} must be set to secondary/peripheral mode with their respective bits in the \register{PxSEL} registers. Finally, the GPIO pins attached to each slave's CS pin need to be set to output a logic high voltage (note that these CS pins are \textit{not} the same as this \peripheral{SPIx} peripheral's own \pin{CSx} pin, which is used when this device is in slave mode).

To initiate a SPI transfer, the master should first clear the SPI status register \register{SPIxSR} by writing 0 to it, and then must assert (pull low) the slave's CS pin. Then, the master may begin transferring as many frames as it wants in the packet. To send a frame, the master must first check if the \register{SPIxTX} register is empty (i.e. ready to accept a new word to queue for transfer) by ensuring that either \bitfield{SPITEIF} is 1 or \bitfield{SPIBUSY} is 0. Then, the master may write the next word to transmit to \register{SPIxTX}. Next, if \bitfield{SPITCIF} is 1, the master should read the value of the \register{SPIxRX} register to read the value the slave transmitted on the prior frame. Then, the master must clear the status register and may either start another frame or end the packet by deasserting (pulling high) the CS wire.

A simplified procedure for producing master mode SPI transfers is shown in Listing \ref{lst:spi-master-transfer} that does not use the TX queue for continuous transfers.

\begin{lstlisting}[style=customC, label=lst:spi-master-transfer, caption={Example Program for SPI Master Transfers}]
/** Includes **/
#include <MemoryMap.h>

/** Function Declarations **/
uint32_t spi_transfer(uint32_t tx_data)

/** Main Function **/
int main()
{
	// Initialize the SPI peripheral
	SPIxCR = SPIEN;

	// Configure the SCKx, MOSIx, and MISOx pins
	SCKx_PxSEL |= SCKx_BIT;
	MOSIx_PxSEL |= MOSIx_BIT;
	MISOx_PxSEL |= MISOx_BIT;

	// Configure the GPIO pin attached to the slave's CS pin as an output high
	PxSEL &= ~(BITy);
	PxOUT |= BITy;
	PxDIR |= BITy;
	
	// Assert the CS pin
	PxOUT &= ~(BITy);

	// Do a single SPI transfer
	uint32_t tx_data = 57;
	uint32_t rx_data;
	
	rx_data = spi_transfer(tx_data);

	// Deassert the CS pin
	PxOUT |= BITy;

	return 0;
}

/** Function Definitions **/
uint32_t spi_transfer(uint32_t tx_data)
{
	// Begin a SPI frame by writing the transmit data to SPIxTX
	SPIxTX = tx_data;

	// Wait for the frame to finish transmitting
	while (SPIxSR & SPIBUSY) {}

	// Return the received data from the slave
	return SPIxRX;
}
\end{lstlisting}



\subsection{Slave Mode Transfer Procedure}

To use the SPI peripheral in slave mode, the user must first initialize the \peripheral{SPIx} peripheral by configuring the \register{SPIxCR} register. Of note, \bitfield{SPIEN} must be 1 and \bitfield{SPISM} must be 1 to enable the SPI peripheral in slave mode. Then, the GPIO pins multiplexed with \pin{SCKx}, \pin{MISOx}, and \pin{MOSIx} must be set to secondary/peripheral mode with their respective bits in the \register{PxSEL} registers. Finally, the \pin{CSx} pin must be set as a high impedance input in GPIO mode to allow the master to assert this \peripheral{SPIx} peripheral.

Immediately after initialization, a word should be written to the \register{SPIxTX} register, which will be transmitted to the master when the master asserts the \pin{CSx} pin and begins sending SCK pulses. Whatever data is contained in \register{SPIxTX} will always be transmitted to the master when the master begins each frame regardless of if the slave has updated the value of \register{SPIxTX}, so care must be taken to always put valid data in the \register{SPIxTX} register before the master begins the next frame. Once the master begins sending a frame, the \bitfield{SPITEIF} flag will become 1, indicating that the next word to transmit to the master should be written to \register{SPIxTX}. At the end of the frame, the \bitfield{SPITCIF} flag will become 1, and the data in \register{SPIxRX} will contain the data received from the master. \register{SPIxRX} should then be read, and the status register \register{SPIxSR} must be cleared. Note that the \register{SPIxRX} register must be read before the end of the next frame, else the data will be overwritten and unrecoverable.

There is no dedicated interrupt indicating when the \pin{CSx} pin is asserted or deasserted, but this functionalty can be achieved using the ordinary GPIO pin interrupts.



\subsection{External SPI Flash Native Memory Read Access}

The \peripheral{SPI0} peripheral, which is attached to the external SPI flash memory containing the program, has a native memory access feature to quickly read data on the SPI flash as if it were addressable by the MCU memory bus. In other words, a pointer can be used in software to read the SPI flash memory as if it were a block of ROM on the address bus. This feature is intended to expand the maximum size of a program from the size of the internal RAM to the size of the SPI flash memory, or 8 MB. Note that using this feature will yield a significant reduction in program execution speed, as each instruction fetched from the SPI flash memory takes either 32 or 64 SCK pulses to read. This feature supports the AT45DB SPI flash memories with either 256- or 264-byte pages, up to the 64 megabit (8 MB) model.

To use the native SPI flash memory access feature, the SPI flash must first be initialized:
\begin{enumerate}
	\item Set up the \peripheral{SPI0} peripheral using master mode, SPIMODE3, 8-bit transfers, an SCK frequency not exceeding 85 MHz, swapped RX bytes, not swapped TX bytes, and the native SPI flash memory access feature disabled (\bitfield{SPIFEN} = 0)
	\item Configure the \pin{SCK0}, \pin{MOSI0}, and \pin{MISO0} pins in secondary/peripheral mode with the appropriate \register{PxSEL} registers
	\item Configure the \pin{CS\_FLASH} pin in GPIO mode, outputting high (deasserted)
	\item Assert (pull low) the CS wire for the SPI flash
	\item Send the 8-bit wake from deep sleep command, \texttt{0xAB}, to the SPI flash
	\item Deassert (pull high) the CS wire and delay for at least 10 ms
	\item Assert (pull low) the CS wire and set the \peripheral{SPI0} peripheral to 32-bit transfers
	\item Send the 32-bit command setting the page length of the SPI flash to 256 bytes: \texttt{0x3D2A7F9A}
	\item Deassert (pull high) the CS wire
	\item Set the \pin{CS\_FLASH} pin to secondary/peripheral mode with the appropriate \register{PxSEL} register
	\item Enable the native SPI flash memory access feature by setting \bitfield{SPIFEN} to 1
\end{enumerate}

After completing the initialization procedure, any 32-bit word stored in the SPI flash may be read by accessing a memory location between \texttt{0x01000000} and \texttt{0x01FFFFFF} on the native MCU address bus. The \peripheral{SPI0} peripheral translates the address from the MCU address bus to the SPI flash using the SPI flash offset register \register{SPI0FOS}. The address accessed on the SPI flash is the 24-bit unsigned addition of the MCU address bus with \register{SPI0FOS}. Note that the addition wraps around 0 if the sum exceeds \texttt{0xFFFFFF}, discarding any bits outside the range of bits 23 down to 0. \register{SPI0FOS} can be used to map a block of data in the SPI flash to a desired MCU address.



\subsection{\peripheral{SPI0} Boot Behavior}

The \peripheral{SPI0} peripheral is used at boot time to load the program from the SPI flash. When the chip is powered on or reset, the state of the \pin{BOOT} pin determines if the chip will fetch and execute the program from the SPI flash or launch the interactive Forth interpreter. If it fetches the program, it uses the \peripheral{SPI0} peripheral to receive the program from the attached SPI flash. As a consequence, the user must guarantee that any other slave devices attached to the \peripheral{SPI0} peripheral will not vie for control over the \pin{MISO0} pin. This error can happen if one of the GPIO pins controls the CS pin of such a slave, as most of the GPIO pins become high impedance inputs upon chip reset. It is therefore possible that the slave's CS wire will settle to an asserted (low) state, causing the slave to attempt to seize control of the \pin{MISO0} pin. Care must be taken to prevent this possibility. One option is to use a start high pin (\pin{SHx}), which resets as an output driving high, to control the slave's CS wire. Another option is to keep the SPI flash as the only slave attached to the \peripheral{SPI0} peripheral.
